% !TeX root = ../se200_referential_regimes.tex

% =============================================================================
\section{The Nine-Regime Lower-Bound Core}
\label{sec:nine}
% =============================================================================

The transformation analysis of Section~\ref{sec:algebra}
yields the following nine core regimes.
Three carrier-basis pairs contribute one regime each.
Three carrier kinds have witnessed split pairs.

The table uses the following short regime labels:

\begin{description}[style=nextline,leftmargin=2.5cm]
  \item[$\OBL$] obligation-bearer identity.
  \item[$\OCC$] occurrence identity.
  \item[$\REC$] record identity.
  \item[$\LOC$] locus-fixed plain-referent identity.
  \item[$\OBJ$] object-fixed plain-referent identity.
  \item[$\SCOPEE$] extension-fixed scope identity.
  \item[$\SCOPES$] structure-fixed scope identity.
  \item[$\RULEC$] content-fixed rule identity.
  \item[$\RULES$] structure-fixed rule identity.
\end{description}

Their classifications over $\Fset$ are collected in Table~\ref{tab:master}.
The table records the core profiles used in the lower-bound proof.
Three same-carrier cells carry the split arguments:
BF for $\LOC/\OBJ$,
AD for $\SCOPEE/\SCOPES$,
and RF for $\RULEC/\RULES$.
Together with carrier disjointness, those cells carry the lower-bound result.
The remaining entries summarize diagnostic behavior of the core profiles.

\begin{table}[ht]
  \centering
  \small
  \setlength{\tabcolsep}{4pt}
  \begin{tabular}{@{}l ccc cc cc cc@{}}
    \toprule
       & OBL  & OCC  & REC  & LOC  & OBJ  & SCOPE-E & SCOPE-S & RULE-C & RULE-S \\
    \midrule
    RE & \NEU & \NEU & \NEU & \NEU & \NEU & \NEU    & \NEU    & \NEU   & \NEU   \\
    AN & \NEU & \NEU & \PRS & \NEU & \NEU & \NEU    & \NEU    & \NEU   & \NEU   \\
    RF & \PRS & \PRS & \PRS & \PRS & \PRS & \PRS    & \BRK    & \PRS   & \BRK   \\
    AD & \PRS & \BRK & \PRS & \PRS & \PRS & \PRS    & \BRK    & \PRS   & \BRK   \\
    RC & \NEU & \NEU & \NEU & \NEU & \NEU & \PRS    & \BRK    & \NEU   & \NEU   \\
    RA & \NEU & \NA  & \NEU & \NA  & \NA  & \NEU    & \BRK    & \NEU   & \NEU   \\
    SU & \BRK & \BRK & \BRK & \BRK & \BRK & \BRK    & \BRK    & \BRK   & \BRK   \\
    BF & \NEU & \NA  & \PRS & \PRS & \BRK & \PRS    & \BRK    & \BRK   & \BRK   \\
    PV & \NEU & \NEU & \PRS & \NEU & \NEU & \NEU    & \NEU    & \NEU   & \NEU   \\
    SE & \NEU & \PRS & \NEU & \PRS & \PRS & \NA     & \NA     & \NA    & \NA    \\
    \bottomrule
  \end{tabular}
  \caption{Classification of each regime profile over the transformation basis $\Fset$.
    \PRS{}: preserves identity;
    \BRK{}: breaks identity;
    \NEU{}: identity-neutral/non-breaking;
    \NA{}: not applicable.
    Each witnessed split pair has a canonical separating operation:
    plain referent at BF, scope at AD, and rule at RF.
    The remaining classifications are diagnostic rather than load-bearing
    for the lower-bound proof.}
  \label{tab:master}
\end{table}

The table should be read as a compact statement of identity behavior.
A \PRS{} entry means that the transformation \emph{preserves} the regime's identity basis.
A \BRK{} entry means that the transformation \emph{breaks} that basis and produces a distinct referent under the regime.
An \NEU{} entry means that the transformation is \emph{identity-neutral} for that regime.
An \NA{} entry means that the transformation is \emph{not applicable} to that regime.

Two columns bracket the profiles.
Re-expression is identity-neutral for every core regime:
a change of representation format preserves the declared identity basis.
Substitution is identity-breaking for every core regime:
replacement of the identity carrier breaks identity under each basis.
These uniform columns serve as classification sanity checks.
They set uniform boundary cases for the profile table.
The lower-bound witnesses are the separating cells for the three split pairs.

The witnessed split pairs carry the lower-bound argument.
$\LOC$ and $\OBJ$ are separated by branching or forking.
$\SCOPEE$ and $\SCOPES$ are separated by aggregation or decomposition.
$\RULEC$ and $\RULES$ are separated by refinement.
The canonical separating operation explains why each split is forced.
The two profiles in a split pair may also differ at other transformations.
For scope,
re-contextualization and reassignment are two such non-canonical differences.
They preserve covered cases and are non-breaking for $\SCOPEE$.
They change applicability context or assignment structure and are breaking for $\SCOPES$.

These regimes are defined by identity basis and transformation behavior.
They refine identity bases within the existing carrier kinds.

\begin{definition}[Separation witness]
  \label{se200.def.SeparationWitness}
  A \emph{separation witness} for identity regimes $\tau$ and $\tau'$ is a pair of representations $(r,s)$ such that
  \[
    r \sim_\tau s
    \qquad\text{and}\qquad
    r \not\sim_{\tau'} s,
  \]
  or conversely.
\end{definition}

\begin{proposition}[Distinct core profiles]
  \label{se200.prop.DistinctCoreProfiles}
  The nine profiles of Table~\ref{tab:master} are pairwise distinct.
\end{proposition}

\begin{proof}
  Same-carrier pairs are separated by the split propositions of Section~\ref{sec:algebra}.
  By Proposition~\ref{se200.prop.PlainReferentUnderdetermination}, $\LOC$ and $\OBJ$ are separated at BF.
  By Proposition~\ref{se200.prop.ScopeUnderdetermination}, $\SCOPEE$ and $\SCOPES$ are separated at AD.
  By Proposition~\ref{se200.prop.RuleUnderdetermination}, $\RULEC$ and $\RULES$ are separated at RF.
  In each case there is a separation witness: a pair identified under one regime and separated under the other.

  Profiles with different carriers govern disjoint sets of representations.
  A representation of an obligation-bearer is never a representation of an occurrence,
  an occurrence never a record,
  a rule never a scope,
  a scope never a plain referent,
  and a record never the thing it records.
  Two regimes whose carriers have disjoint representation domains induce
  identity relations that differ on those domains,
  and are therefore distinct by Definition~\ref{se200.def.Collapse}.
  Identifying such profiles would collapse the reference inventory of
  Section~\ref{sec:inventory} and force one carrier to stand in for another.

  Therefore every pair of core regimes in Table~\ref{tab:master}
  is separated either by a same-carrier transformation witness
  or by carrier difference.
  Appendix~\ref{app:sep} records the separation basis for all $\binom{9}{2}=36$ pairs.
\end{proof}
